% 第三章 系统需求分析
% 建议文件路径：chapters/3_chapter3.tex
% 写作要求：基于论文目录大纲第三章扩展；图表均在正文中引出和分析；每个小节尽量有文献引用。
% 风格说明：学术表达约 70%，自然化说明约 30%。

\chapter{系统需求分析}
\label{chap:requirements-analysis}

系统需求分析是系统设计与实现之前的重要环节。对于在线视频分享平台而言，需求分析不能只停留在“能播放视频”这一层面，而需要从用户角色、业务流程、功能边界、内容审核、性能要求和系统可行性等方面进行综合分析。只有把需求边界梳理清楚，后续的总体架构设计、数据库设计和功能实现才不会偏离系统目标。

根据论文目录大纲中第三章的安排，本章主要围绕系统总体需求、用户端与创作中心需求、管理端与 AI 审核需求、非功能需求以及可行性分析展开。与前一章的相关技术介绍不同，本章更关注“系统需要做什么”和“为什么需要这样设计”。为了让分析过程更清晰，本章加入了角色用例图、业务流程图、状态流转图和多张需求表，并在正文中对图表进行说明。

\section{系统总体需求分析}
\label{sec:overall-requirements}

\subsection{需求分析目标}
\label{subsec:req-analysis-goal}

本课题旨在设计并实现一个基于微服务架构的在线视频分享平台。系统需要围绕视频内容的生产、处理、分发、互动和审核形成完整业务闭环。普通用户可以浏览、搜索和播放视频，并参与评论、弹幕、点赞、收藏、投币等互动；内容创作者可以上传视频、编辑稿件和管理个人作品；管理员可以在后台进行分类维护、视频审核和内容治理；AI 审核服务则作为独立能力，对上传视频进行自动化风险识别，为人工审核提供辅助依据。

从软件工程角度看，需求分析需要先明确系统参与者、核心业务流程和质量目标，再进一步细化功能需求和非功能需求。软件架构实践中也强调，需求分析不仅要关注功能本身，还要关注系统质量属性，例如性能、可用性、安全性和可维护性\cite{bass2021softwarearchitecture}。因此，本章在分析功能需求的同时，也会对系统的非功能需求进行说明。

本系统的总体目标不是简单实现一个视频播放器，而是构建一个从视频投稿、视频处理、内容审核到正式发布、用户播放和社区互动的完整平台。也就是说，视频内容从创作者提交开始，需要经过上传、转码、审核、发布等步骤，最后才能被普通用户浏览和播放。这个过程涉及多个角色和多个业务环节，需要在需求分析阶段提前梳理清楚。

\subsection{系统参与角色划分}
\label{subsec:system-actors}

系统参与角色是需求分析中的基础内容。不同角色拥有不同操作权限，也对应不同功能需求。对于在线视频分享平台来说，如果不区分游客、注册用户、创作者、管理员和 AI 审核服务，很容易把所有功能混在一起，导致后续权限设计和业务流程不清晰。

本系统主要参与角色包括游客、注册用户、内容创作者、管理员和 AI 审核服务。其中，游客只能进行基础浏览和搜索；注册用户在登录后可以参与互动；内容创作者可以上传和管理视频；管理员负责后台管理和人工审核；AI 审核服务负责处理视频审核任务并返回审核结果。角色划分的目的不是增加复杂度，而是让每类用户的功能边界更加明确。

表~\ref{tab:system-actors} 展示了系统主要角色及其需求说明。通过该表可以看到，普通用户侧重内容消费，创作者侧重内容生产，管理员侧重内容治理，而 AI 审核服务则承担自动化辅助分析任务。

\begin{table}[htbp]
  \centering
  \caption{系统主要角色及功能需求}
  \label{tab:system-actors}
  \begin{tabular}{p{0.20\textwidth}p{0.36\textwidth}p{0.34\textwidth}}
    \hline
    \textbf{角色} & \textbf{主要功能需求} & \textbf{说明} \\
    \hline
    游客 & 浏览视频、搜索视频、查看公开视频信息 & 未登录用户可进行基础访问 \\
    注册用户 & 登录、播放视频、评论、弹幕、点赞、收藏、投币、关注 & 登录后可参与平台互动 \\
    内容创作者 & 上传视频、编辑投稿、查看稿件状态、管理作品 & 负责视频内容创作与发布 \\
    管理员 & 后台登录、分类管理、视频审核、查看 AI 审核结果 & 负责平台内容管理与审核 \\
    AI 审核服务 & 消费审核任务、分析视频、返回审核结果 & 为管理员审核提供辅助依据 \\
    \hline
  \end{tabular}
\end{table}

由表~\ref{tab:system-actors} 可以看出，系统并不是只有“用户”和“管理员”两个简单角色。内容创作者虽然属于注册用户，但其操作范围更偏向视频投稿和稿件管理；AI 审核服务虽然不是传统意义上的人工角色，但它会参与视频审核流程，因此在需求分析中也需要单独列出。

\subsection{系统总体业务流程}
\label{subsec:overall-business-flow}

在线视频分享平台的核心流程可以概括为“内容生产、内容处理、内容审核、内容发布、内容消费”。内容创作者首先提交视频，系统对视频进行上传和转码处理；转码完成后，视频进入审核阶段；审核通过后，视频才能进入正式展示和播放流程。这个流程体现了视频平台对内容安全和资源处理的基本要求。

为了更直观地展示系统总体业务过程，图~\ref{fig:overall-business-flow} 给出了系统核心业务流程示意图。该图从内容创作者上传视频开始，到用户播放和互动结束，展示了平台业务闭环。

\begin{figure}[htbp]
  \centering
  \fbox{%
    \begin{minipage}{0.90\textwidth}
      \centering
      \textbf{在线视频分享平台总体业务流程}\\[0.8em]
      创作者提交视频 $\rightarrow$ 视频上传与投稿保存\\[0.4em]
      $\Downarrow$\\[0.4em]
      视频转码与资源处理 $\rightarrow$ AI 初步审核\\[0.4em]
      $\Downarrow$\\[0.4em]
      管理员人工复核 $\rightarrow$ 审核通过后正式发布\\[0.4em]
      $\Downarrow$\\[0.4em]
      用户浏览、搜索、播放视频 $\rightarrow$ 评论、弹幕、点赞等互动
    \end{minipage}
  }
  \caption{系统总体业务流程示意图}
  \label{fig:overall-business-flow}
\end{figure}

由图~\ref{fig:overall-business-flow} 可以看出，本系统的业务流程并不是一次性完成的同步流程，而是包含上传、处理、审核、发布和互动多个阶段。尤其是视频转码和 AI 审核属于耗时任务，因此需要通过状态流转和异步处理保证流程稳定。数据密集型应用设计中也强调，复杂系统需要通过任务状态、可靠处理和异常恢复来保证长期可维护性\cite{kleppmann2017ddia}。

\subsection{系统用例关系分析}
\label{subsec:system-usecase-analysis}

系统用例分析可以帮助说明不同角色与功能之间的关系。通过用例分析，可以较清楚地看到某类角色能执行哪些操作，以及不同用例之间是否存在依赖关系。对于本系统来说，用例分析重点关注普通用户、内容创作者、管理员和 AI 审核服务之间的功能边界。

图~\ref{fig:system-usecase} 展示了系统主要角色与核心功能之间的关系。

\begin{figure}[htbp]
  \centering
  \fbox{%
    \begin{minipage}{0.90\textwidth}
      \centering
      \textbf{系统角色与功能用例关系}\\[0.8em]
      游客：浏览视频、搜索视频\\[0.4em]
      注册用户：播放视频、评论、弹幕、点赞、收藏、投币、关注\\[0.4em]
      内容创作者：上传视频、提交投稿、查看稿件状态、管理作品\\[0.4em]
      管理员：后台登录、分类管理、视频审核、查看 AI 审核结果\\[0.4em]
      AI 审核服务：消费审核任务、分析视频内容、返回审核结果
    \end{minipage}
  }
  \caption{系统角色与功能用例关系示意图}
  \label{fig:system-usecase}
\end{figure}

由图~\ref{fig:system-usecase} 可以看出，游客、注册用户和内容创作者之间存在权限递进关系。游客只能访问公开视频内容；注册用户在此基础上可以参与互动；内容创作者进一步拥有投稿和作品管理能力。管理员与普通用户侧的功能不同，主要承担内容管理和审核职责。AI 审核服务则通过自动化分析参与审核流程，但不直接决定视频最终发布结果。

\section{用户端功能需求分析}
\label{sec:user-side-requirements}

\subsection{游客访问需求}
\label{subsec:visitor-requirements}

游客是未登录状态下访问平台的用户。对于一个公开视频分享平台来说，游客访问能力是基础需求，因为平台需要允许未登录用户浏览公开视频内容，以降低用户进入平台的门槛。游客主要可以查看首页视频列表、根据分类浏览视频、通过关键词搜索视频，以及进入公开视频详情页查看基础信息。

不过，游客的操作范围需要受到限制。涉及个人身份和用户行为的数据操作，例如发表评论、发送弹幕、点赞、收藏、投币和关注等，都应要求用户登录后才能进行。这样的设计既能保证用户行为可追踪，也能避免匿名操作带来的数据混乱和安全问题。权限控制通常需要结合认证与授权机制进行设计，Spring Security 文档也对认证和访问控制能力进行了说明\cite{springsecuritydocs}。

游客访问需求可以概括为“可浏览、可搜索、少交互”。也就是说，游客能够了解平台内容，但不能对平台内容产生需要身份绑定的行为。这样既保证了开放访问体验，也为用户登录和注册提供了自然引导。

\subsection{用户账号与身份认证需求}
\label{subsec:user-auth-requirements}

注册用户是平台中最重要的普通用户角色。系统需要支持用户注册、登录、自动登录、退出登录和登录状态校验等功能。用户登录后，系统应能够识别当前用户身份，并根据身份允许其进行评论、弹幕、点赞、收藏、投币和关注等操作。

用户认证需求不仅包括“能登录”，还包括登录状态的安全维护。系统需要区分普通用户和管理员身份，不能让普通用户访问后台管理接口，也不能让未登录用户执行需要身份记录的操作。登录态可以通过 Token、缓存和后端校验等方式维护，其中缓存系统可用于保存短期登录状态和验证码等信息\cite{redisdocs}。

从用户体验角度看，账号功能应尽量简单直观。用户在未登录时进行需要身份的操作，系统应给出明确提示，而不是直接报错。用户登录后，页面也应能及时显示用户信息和可用操作入口。

\subsection{视频浏览、搜索与播放需求}
\label{subsec:video-browse-search-play}

视频浏览和播放是用户端的核心需求。用户进入平台后，需要能够浏览视频列表，查看视频封面、标题、作者、播放量、弹幕数和发布时间等信息。点击视频后，用户应进入播放页，查看视频详情并进行播放。

搜索功能同样重要。对于内容类平台来说，用户不一定只通过首页浏览内容，也可能通过关键词查找特定视频。搜索需求需要支持标题、简介、标签等字段的匹配，并能够根据相关性、发布时间或热度进行排序。Elasticsearch 等全文检索技术可以为复杂关键词检索提供支持\cite{elasticsearchdocs}。

视频播放方面，由于上传视频通常需要经过转码和切片处理，因此播放页应支持加载流媒体播放资源。HLS 协议通过播放列表和媒体分片组织视频资源，适合网页端视频在线播放场景\cite{rfc8216hls}。从用户角度看，播放过程应尽量稳定，播放失败时也需要给出明确提示。

表~\ref{tab:user-main-functions} 总结了用户端主要功能需求。该表可以看出，用户端功能不仅包括基础浏览，还包括搜索、播放和互动等多个层面。

\begin{table}[htbp]
  \centering
  \caption{用户端主要功能需求}
  \label{tab:user-main-functions}
  \begin{tabular}{p{0.22\textwidth}p{0.32\textwidth}p{0.36\textwidth}}
    \hline
    \textbf{功能类别} & \textbf{功能点} & \textbf{需求说明} \\
    \hline
    基础浏览 & 首页视频列表、分类筛选 & 用户能够查看公开视频内容 \\
    搜索检索 & 关键词搜索、排序、热词 & 用户能够根据标题、简介、标签查找视频 \\
    视频播放 & 视频详情、分 P 切换、在线播放 & 用户能够稳定播放已发布视频 \\
    用户互动 & 评论、弹幕、点赞、收藏、投币、关注 & 登录用户能够参与社区互动 \\
    账号功能 & 注册、登录、退出、自动登录 & 支持用户身份识别和状态维护 \\
    \hline
  \end{tabular}
\end{table}

由表~\ref{tab:user-main-functions} 可以看出，用户端需求具有明显的层次性。浏览和搜索属于开放性需求，播放和详情查看属于内容消费需求，评论弹幕和点赞收藏则属于登录后的互动需求。后续系统设计时，需要根据这些层次合理划分接口权限和页面功能。

\subsection{社区互动需求}
\label{subsec:interaction-requirements}

社区互动是视频分享平台区别于普通视频播放工具的重要特点。用户不仅要能观看视频，还需要能表达自己的态度和观点。评论、弹幕、点赞、收藏、投币和关注等功能，可以增强用户参与感，也能让平台内容形成一定社区氛围。

评论功能需要支持用户在视频下发表观点，也可以支持回复和评论点赞等扩展操作。弹幕功能则与视频播放时间点紧密相关，用户发送弹幕后，弹幕应能在对应播放时间附近展示。弹幕数据一般需要包含视频编号、分 P 文件编号、播放时间、弹幕内容、颜色和模式等信息。

点赞、收藏、投币和关注属于用户行为记录。系统需要避免同一用户重复执行不合理操作，例如重复点赞或重复收藏。对于这类行为，系统不仅要保存操作结果，还要能在用户再次进入页面时展示当前状态。用户行为数据属于高频交互数据，在设计时需要考虑访问效率和数据一致性。

\section{创作中心功能需求分析}
\label{sec:creator-center-requirements}

\subsection{视频投稿需求}
\label{subsec:video-submit-requirements}

创作中心面向内容创作者，主要承担视频投稿和作品管理功能。创作者需要能够填写视频标题、选择视频分类、上传封面、填写标签和简介，并设置投稿相关选项。投稿功能的目标是让创作者可以把视频内容提交给平台，而不是简单上传一个文件。

视频投稿应支持多分 P 或多文件场景。对于一些较长视频或系列视频，创作者可能需要上传多个视频片段。系统需要能够保存每个分 P 的标题、文件信息和排序关系。这样用户在播放页才能进行分 P 切换，管理员在审核时也能分别查看不同分 P 内容。

投稿完成后，视频不应立即公开发布，而应进入处理和审核流程。这样可以避免未经转码或未经审核的视频直接出现在前台。内容平台治理研究也强调，用户生成内容平台需要在发布流程中加入必要的审核和治理机制\cite{gorwa2020algorithmic}。

\subsection{视频上传与文件处理需求}
\label{subsec:upload-file-processing-requirements}

视频文件通常体积较大，直接一次性上传容易受到网络波动影响。因此，系统需要支持大文件分片上传。创作者选择视频文件后，前端可以将文件拆分为多个分片，再逐个上传到服务端。服务端接收分片后保存临时文件，并在所有分片上传完成后进行合并。

上传完成后，系统还需要对视频进行转码处理。用户上传的视频可能来自不同设备，格式和编码方式不完全一致。如果不进行统一处理，可能导致前台播放兼容性差。FFmpeg 可用于完成视频转码、格式转换和切片等处理\cite{ffmpegdocs}。

从需求角度看，视频上传和文件处理过程不能让用户长时间等待。因此，用户提交投稿后，系统应允许后台继续处理视频，并通过稿件状态展示处理进度。这样对用户更友好，也更符合大文件处理的实际情况。

\subsection{稿件状态管理需求}
\label{subsec:post-status-requirements}

稿件状态管理是创作中心的重要需求。创作者提交视频后，需要知道视频当前处于什么阶段，例如上传完成、转码中、转码失败、待审核、审核通过或审核不通过。如果没有状态展示，创作者很难判断投稿是否成功，也不知道下一步是否需要修改。

稿件状态应能反映视频生命周期中的关键节点。比如视频刚提交时可以处于转码中，转码完成后进入待审核，管理员审核通过后进入正式发布状态，审核不通过则应显示驳回状态和原因。状态设计应尽量清晰，避免出现用户看不懂的内部状态名称。

从系统设计角度看，状态管理也有助于后台处理异常情况。比如转码失败时，系统可以记录失败原因；AI 审核失败时，可以转入人工复核。微服务模式中也强调，长流程任务需要通过明确状态和事件处理来保证业务过程可追踪\cite{richardson2018microservices}。

\subsection{创作者作品管理需求}
\label{subsec:creator-work-management}

创作者除了提交视频，还需要管理自己已经提交或发布的作品。作品管理功能应支持查看稿件列表、查看视频状态、编辑未发布稿件、删除不需要的稿件，以及查看作品的基础统计信息。这样创作者才能较完整地管理自己的内容。

作品管理和普通用户浏览不同，它更关注“我上传了什么”和“现在处理到哪一步”。因此，作品管理页面应突出视频状态、投稿时间、审核结果和基础数据，而不是只展示视频封面和播放量。

表~\ref{tab:creator-functions} 对创作中心主要功能需求进行了整理。可以看出，创作中心功能从投稿填写开始，贯穿上传、转码、审核状态查看和作品管理多个环节。

\begin{table}[htbp]
  \centering
  \caption{创作中心功能需求}
  \label{tab:creator-functions}
  \begin{tabular}{p{0.22\textwidth}p{0.32\textwidth}p{0.36\textwidth}}
    \hline
    \textbf{功能类别} & \textbf{功能点} & \textbf{需求说明} \\
    \hline
    投稿信息 & 标题、分类、标签、简介、封面 & 保存视频基础内容信息 \\
    视频上传 & 分片上传、多文件上传 & 支持大文件稳定上传 \\
    文件处理 & 分片合并、视频转码、播放资源生成 & 使视频满足在线播放要求 \\
    状态查看 & 转码中、待审核、审核通过、审核失败 & 创作者可查看稿件处理进度 \\
    作品管理 & 稿件列表、编辑、删除、统计查看 & 创作者可管理个人投稿内容 \\
    \hline
  \end{tabular}
\end{table}

由表~\ref{tab:creator-functions} 可知，创作中心不是单一上传页面，而是一个围绕视频投稿生命周期展开的功能模块。它既要支持内容提交，也要支持处理进度反馈和作品管理。

\section{管理端与 AI 审核需求分析}
\label{sec:admin-ai-requirements}

\subsection{管理员认证与权限需求}
\label{subsec:admin-auth-requirements}

管理端主要面向平台管理员，管理员需要通过后台登录后才能进入管理系统。由于管理端涉及分类维护、视频审核和内容治理等敏感操作，因此管理员认证和权限控制是后台功能的基础。

管理员登录应与普通用户登录区分开来。普通用户不能直接访问后台接口，管理员登录态也不应与普通用户登录态混用。后台接口需要进行身份校验，未登录或权限不足时应拒绝访问。认证与授权机制是系统安全的重要组成部分，Spring Security 对相关能力进行了系统说明\cite{springsecuritydocs}。

从需求角度看，管理端权限控制不需要一开始就设计得特别复杂，但至少应保证后台入口不可被普通用户随意访问。随着系统后续扩展，可以进一步增加管理员角色分级、操作日志和权限细分等功能。

\subsection{分类管理需求}
\label{subsec:category-management-requirements}

分类管理是后台管理中的基础功能。视频分类会影响用户投稿时的分类选择，也会影响前台视频展示、筛选和检索。因此，系统需要支持管理员维护视频分类数据，包括新增分类、修改分类、删除分类、调整排序和启用停用等操作。

分类结构可以分为一级分类和二级分类。一级分类用于表示较大的内容方向，二级分类用于进一步细分内容领域。这样的分类方式既方便用户浏览，也方便创作者选择合适分类。

分类管理需求看似简单，但它会影响视频投稿、首页展示和搜索筛选等多个功能。因此，分类数据需要保持稳定，删除或修改分类时也应考虑是否影响已有视频内容。

\subsection{视频人工审核需求}
\label{subsec:manual-audit-requirements}

视频人工审核是管理端的核心需求。创作者提交视频后，视频不能直接进入正式发布状态，而需要经过转码和审核流程。管理员需要能够查看待审核视频列表，进入审核详情页，观看视频内容，查看投稿信息和 AI 审核结果，并做出通过或驳回决定。

人工审核的价值在于结合平台规则和实际语境做最终判断。AI 审核可以提供风险提示，但模型可能误判，因此最终发布决策仍应由管理员确认。人机交互研究中也指出，AI 系统应为人工决策提供可理解、可纠错的支持，而不是完全替代人工判断\cite{amershi2019guidelines}。

审核通过后，系统应将视频从投稿态转为正式发布态，并允许用户在前台浏览和播放；审核驳回后，系统应保留驳回状态和必要原因，方便创作者了解问题。这样可以形成比较清晰的内容审核闭环。

\subsection{AI 审核任务需求}
\label{subsec:ai-task-requirements}

AI 审核服务的需求主要体现在审核任务自动触发、视频内容多模态分析和审核结果结构化返回三个方面。当视频转码完成后，后端应创建审核任务，并通过消息队列通知 AI 服务进行处理。消息队列适合这类耗时任务，因为它能够让业务系统和 AI 服务解耦\cite{rabbitmqdocs}。

AI 服务收到任务后，需要对视频进行分析。分析内容可以包括语音识别、声音事件识别、关键帧抽取和多模态语义判断。语音识别可以将视频中的人声转换为文本，声音事件识别可以发现异常声音，多模态模型则可以结合画面和文本进行综合判断。

AI 审核结果需要结构化返回，不能只返回一个简单的“通过”或“不通过”。更合理的结果应包括审核建议、风险等级、风险标签、命中时间段、命中原因和关键帧截图等信息。这样管理员在审核时才能快速定位风险点。

\subsection{AI 审核结果展示需求}
\label{subsec:ai-result-display-requirements}

AI 审核结果的价值不只在于后台是否收到了结果，还在于管理员能否看懂结果。如果 AI 只返回一段模糊文字，管理员仍然需要花大量时间重新判断。因此，系统需要在审核详情页中清晰展示 AI 审核建议、风险等级、风险标签和命中片段。

语音识别、声音事件识别和多模态理解分别对应不同类型的风险信息。Whisper 可用于语音转写，AudioSet 和相关模型可用于声音事件识别，Qwen-VL 等视觉语言模型可用于图像与文本联合理解\cite{radford2023whisper}。这些结果在后台展示时，应尽量转化为管理员可以直接理解的信息。

表~\ref{tab:admin-ai-functions} 对管理端与 AI 审核相关需求进行了总结。该表可以看出，管理端并不是单纯执行人工审核，而是需要结合分类管理、视频审核、AI 结果查看和最终发布控制等功能。

\begin{table}[htbp]
  \centering
  \caption{管理端与 AI 审核功能需求}
  \label{tab:admin-ai-functions}
  \begin{tabular}{p{0.24\textwidth}p{0.32\textwidth}p{0.34\textwidth}}
    \hline
    \textbf{功能类别} & \textbf{功能点} & \textbf{需求说明} \\
    \hline
    管理员认证 & 后台登录、退出、权限校验 & 保证后台功能只对管理员开放 \\
    分类管理 & 分类新增、编辑、删除、排序 & 维护平台内容分类体系 \\
    视频审核 & 审核列表、审核详情、通过、驳回 & 控制视频是否正式发布 \\
    AI 审核任务 & 自动创建任务、投递请求、接收结果 & 辅助完成视频内容风险识别 \\
    AI 结果展示 & 风险等级、风险标签、命中片段、原因说明 & 帮助管理员快速定位风险内容 \\
    \hline
  \end{tabular}
\end{table}

由表~\ref{tab:admin-ai-functions} 可知，AI 审核并不是替代管理端审核，而是增强审核能力。管理员仍然是最终审核者，AI 服务主要负责提前分析视频内容并提供风险提示。

\subsection{视频状态流转需求}
\label{subsec:video-status-flow-requirements}

为了保证审核流程清晰，系统中视频状态需要明确流转。投稿视频从用户提交开始，通常会依次经历转码中、转码失败、待审核、审核通过或审核不通过等状态。AI 审核结果不直接决定视频最终发布，而是作为管理员人工审核的重要参考。

图~\ref{fig:video-status-flow} 展示了视频投稿与审核状态流转过程。通过该图可以看出，视频正式发布之前需要经过转码和审核两个关键阶段。

\begin{figure}[htbp]
  \centering
  \fbox{%
    \begin{minipage}{0.90\textwidth}
      \centering
      \textbf{视频投稿与审核状态流转}\\[0.8em]
      用户提交投稿 $\rightarrow$ 转码中\\[0.4em]
      转码中 $\rightarrow$ 转码失败 或 待审核\\[0.4em]
      待审核 $\rightarrow$ 创建 AI 审核任务 $\rightarrow$ 返回 AI 审核结果\\[0.4em]
      管理员人工审核 $\rightarrow$ 审核通过并正式发布 或 审核不通过
    \end{minipage}
  }
  \caption{视频投稿与审核状态流转图}
  \label{fig:video-status-flow}
\end{figure}

由图~\ref{fig:video-status-flow} 可以看出，状态流转是视频发布流程中的核心控制手段。只有当转码完成且审核通过后，视频才能进入正式发布状态。这样既可以避免无法播放的视频进入前台，也可以减少未审核内容直接展示带来的风险。

\section{非功能需求分析}
\label{sec:non-functional-requirements}

\subsection{性能需求}
\label{subsec:performance-requirements}

性能需求主要关注系统在用户访问时是否具有较好的响应速度。对于在线视频分享平台来说，常见页面和接口应保持流畅，例如登录、视频列表加载、视频详情查询、评论加载和搜索查询等。用户不一定能感知系统内部架构，但会直接感受到页面是否卡顿、搜索是否慢、视频是否能正常播放。

对于视频上传、视频转码和 AI 审核等耗时操作，系统不应要求用户一直等待处理完成，而应采用异步处理方式。ISO/IEC 25010 软件质量模型将性能效率作为软件质量的重要方面，其中包括时间特性、资源利用和容量等内容\cite{iso25010}。

因此，性能需求可以分为两类：一类是普通访问接口需要尽量快速响应；另一类是耗时任务需要进入后台处理，并通过状态反馈给用户。这样既能保证页面体验，也能让复杂任务有足够时间完成。

\subsection{可扩展性需求}
\label{subsec:scalability-requirements}

可扩展性需求关注系统后续增加功能或提升处理能力是否方便。在线视频分享平台可能随着用户数量、视频数量和互动数据增长而面临更高压力。如果系统结构过于集中，后续扩展就会比较困难。

微服务架构的一个重要优势是可以按业务模块进行扩展。例如，当视频资源处理压力较大时，可以优先扩展资源处理相关服务；当互动数据访问频繁时，可以优化或扩展互动模块。微服务模式中也强调，服务应围绕业务能力拆分，以便独立部署和扩展\cite{richardson2018microservices}。

不过，可扩展性并不意味着一开始就把所有功能拆得很细。过度拆分会增加运维和通信成本。因此，本系统的可扩展性需求应理解为：当前结构能够支撑后续扩展，而不是在毕业设计阶段追求过度复杂的部署方式。

\subsection{安全性需求}
\label{subsec:security-requirements}

安全性需求主要体现在用户身份认证、管理员权限控制和内部接口保护等方面。普通用户和管理员应使用不同的登录态标识，后台管理接口需要进行管理员身份校验，普通用户不应直接访问后台或内部接口。

系统还需要对验证码、Token、上传文件和审核接口进行必要保护。比如用户登录注册需要验证码辅助校验，视频上传需要限制文件类型和大小，后台审核接口需要检查管理员身份。Spring Security 提供了认证、授权和访问控制方面的能力说明，可作为权限设计的技术参考\cite{springsecuritydocs}。

从需求角度看，安全性不是最后附加的功能，而应贯穿登录、上传、审核和管理等关键流程。尤其是后台审核功能，如果权限控制不当，可能会直接影响平台内容安全。

\subsection{可靠性与一致性需求}
\label{subsec:reliability-consistency-requirements}

可靠性需求关注系统在异常情况下是否仍能保持基本可用。例如视频转码失败、AI 审核失败、消息重复投递或服务短暂不可用时，系统不应直接崩溃，也不应让业务状态混乱。可靠性设计通常需要结合状态记录、异常处理和失败恢复进行考虑。

一致性需求则主要关注投稿数据、视频文件状态、审核状态和正式发布数据之间是否保持逻辑一致。例如，如果视频还没有审核通过，就不应该进入正式视频列表；如果视频转码失败，也不应该允许用户播放。分布式系统中，一致性往往需要结合事务、状态流转和最终一致性策略处理\cite{kleppmann2017ddia}。

在本系统中，视频发布流程较长，因此状态记录尤为重要。只要状态设计清楚，即使某个环节出现失败，也可以通过后台查看和人工处理进行恢复。

\subsection{可维护性需求}
\label{subsec:maintainability-requirements}

可维护性需求关注系统是否便于后续修改、扩展和排查问题。毕业设计系统虽然规模有限，但仍然需要保持合理的目录结构、模块边界和接口规范。否则随着功能增加，代码会很快变得难以维护。

可维护性可以通过多种方式提升，例如前后端分离、服务职责划分、公共工具封装、统一异常处理、接口文档化和日志记录等。软件架构实践中也强调，可维护性是系统长期演进的重要质量属性\cite{bass2021softwarearchitecture}。

表~\ref{tab:non-functional-requirements} 对系统主要非功能需求进行了总结。该表可以看出，非功能需求不是单独存在的，而是和系统功能实现紧密相关。

\begin{table}[htbp]
  \centering
  \caption{系统主要非功能需求}
  \label{tab:non-functional-requirements}
  \begin{tabular}{p{0.20\textwidth}p{0.36\textwidth}p{0.34\textwidth}}
    \hline
    \textbf{非功能需求} & \textbf{需求说明} & \textbf{设计关注点} \\
    \hline
    性能需求 & 常规页面和接口响应应保持流畅 & 缓存、搜索、异步任务处理 \\
    可扩展性 & 系统应便于后续增加功能和扩展服务 & 模块拆分、服务边界、接口规范 \\
    安全性 & 区分普通用户和管理员权限，保护内部接口 & 身份认证、权限校验、上传限制 \\
    可靠性 & 异常情况下核心流程应尽量不中断 & 状态记录、失败原因、人工处理 \\
    数据一致性 & 投稿、转码、审核和发布状态应保持一致 & 事务控制、状态流转、结果回写 \\
    可维护性 & 系统结构清晰，便于后续修改和扩展 & 分层设计、公共模块、日志和文档 \\
    \hline
  \end{tabular}
\end{table}

由表~\ref{tab:non-functional-requirements} 可知，非功能需求实际上决定了系统能否稳定运行。功能实现只能说明系统“能不能用”，而非功能需求则影响系统“好不好用、稳不稳定、能不能继续扩展”。

\section{系统可行性分析}
\label{sec:feasibility-analysis}

\subsection{技术可行性}
\label{subsec:technical-feasibility}

从技术角度看，本系统所涉及的关键技术均较为成熟。后端服务可以基于 Java 和微服务框架实现，数据存储可以使用关系型数据库和缓存系统，视频处理可以使用成熟的音视频工具，AI 审核也可以通过已有模型和独立服务完成。整体技术路线具有可实现性。

当然，技术可行并不代表实现过程没有难点。系统的难点主要集中在视频分片上传、视频转码、审核状态流转、AI 审核结果回写和前后端联调等方面。这些问题需要在实现阶段通过模块化设计、状态字段和接口联调逐步解决。

从毕业设计角度看，本系统选择的技术既能体现工程完整性，也不会脱离实际开发能力范围。视频处理和 AI 审核具有一定挑战性，但可以通过合理拆分功能和逐步实现来完成。

\subsection{经济可行性}
\label{subsec:economic-feasibility}

从经济角度看，本系统主要基于开源技术和本地开发环境实现，不需要购买额外商业软件授权。Java、MySQL、Redis、Elasticsearch、RabbitMQ、FFmpeg 和常见 AI 模型工具均可以在普通开发环境中进行学习和实验。

系统开发过程中的主要成本是时间成本和硬件资源成本。普通业务服务和前端页面可以在常规电脑上运行，AI 审核部分如果使用较大的多模态模型，则可能对 GPU 和显存有一定要求。因此，在毕业设计阶段可以根据实际硬件条件选择模型规模或降低推理复杂度。

总体来看，本系统在经济上是可行的。即使没有高性能服务器，也可以先完成核心业务流程，再根据条件逐步完善 AI 审核能力。

\subsection{操作可行性}
\label{subsec:operational-feasibility}

从操作角度看，系统面向的用户主要包括普通用户、内容创作者和管理员。普通用户通过浏览器访问平台，完成浏览、搜索、播放和互动；内容创作者通过创作中心完成投稿；管理员通过后台页面完成分类管理和视频审核。这些操作都属于 Web 系统中较常见的交互方式。

为了保证操作可行性，系统界面应尽量保持清晰。普通用户不需要理解视频转码和审核流程，只需要看到视频是否能播放；创作者需要看到投稿状态；管理员需要看到待审核视频和 AI 风险结果。只要不同角色看到的信息与其任务相匹配，系统使用难度就不会太高。

表~\ref{tab:feasibility-summary} 对系统可行性进行了总结。可以看出，本系统在技术、经济和操作三个方面均具备实现条件。

\begin{table}[htbp]
  \centering
  \caption{系统可行性分析汇总}
  \label{tab:feasibility-summary}
  \begin{tabular}{p{0.22\textwidth}p{0.42\textwidth}p{0.26\textwidth}}
    \hline
    \textbf{可行性类别} & \textbf{分析内容} & \textbf{结论} \\
    \hline
    技术可行性 & 主要技术成熟，系统可通过分阶段实现完成 & 可行 \\
    经济可行性 & 基于开源技术和本地环境开发，成本较低 & 可行 \\
    操作可行性 & 用户和管理员均通过 Web 页面操作，交互方式直观 & 可行 \\
    \hline
  \end{tabular}
\end{table}

由表~\ref{tab:feasibility-summary} 可以看出，本系统虽然涉及视频处理和 AI 审核等复杂内容，但整体仍可以通过开源技术和分阶段开发完成。对于毕业设计而言，重点在于保证核心流程完整，而不是一开始就追求商业平台级别的性能和规模。

\section{本章小结}
\label{sec:chapter3-summary}

本章对在线视频分享平台进行了系统需求分析。首先，从总体需求出发，明确了系统需要围绕视频内容生产、处理、审核、发布、播放和互动形成完整闭环，并通过表~\ref{tab:system-actors}、图~\ref{fig:overall-business-flow} 和图~\ref{fig:system-usecase} 分析了系统角色、业务流程和用例关系。

随后，本章分别分析了用户端、创作中心、管理端与 AI 审核相关功能需求。用户端主要关注视频浏览、搜索、播放和社区互动；创作中心主要关注视频投稿、分片上传、稿件状态和作品管理；管理端主要关注管理员认证、分类管理、视频审核和 AI 审核结果展示。图~\ref{fig:video-status-flow} 进一步说明了视频从投稿到正式发布的状态流转过程。

最后，本章从性能、可扩展性、安全性、可靠性、一致性和可维护性等方面分析了系统非功能需求，并从技术、经济和操作三个方面论证了系统可行性。通过本章分析，可以为第四章的系统总体设计提供较清晰的需求基础。
